\documentclass[11pt]{article}
\usepackage[margin=1in]{geometry}
\usepackage{amsmath,amssymb,amsthm}
\usepackage{array}
\usepackage{parskip}
\usepackage{booktabs}
\usepackage{ctex}

\newtheorem{theorem}{Theorem}
\newtheorem{lemma}[theorem]{Lemma}
\newtheorem{corollary}[theorem]{Corollary}
\newtheorem{proposition}[theorem]{Proposition}
\theoremstyle{definition}
\newtheorem{definition}[theorem]{Definition}
\theoremstyle{remark}
\newtheorem{remark}[theorem]{Remark}

\newcommand{\F}{\mathbb{F}}
\newcommand{\PG}{\mathrm{PG}}
\newcommand{\ov}{\mathcal{O}}

\title{Disproof of conjecture \texttt{00000001006}}
\author{earthking11}
\date{2026-09-13}

\begin{document}
\maketitle

\section*{Verdict: FALSE}

We disprove conjecture \texttt{00000001006} as stated. The conjecture asserts
that an $m$-ovoid of the parabolic quadric $Q(4,q)$ can exist only when
$m \mid (q+1)$. At the smallest case $q=2$, $m=2$ the claim fails: $Q(4,2)$ has
an ovoid (a $1$-ovoid), and the complement of that ovoid meets each of the
$15$ generating lines in exactly $2$ points, so it is a $2$-ovoid. But
$2 \nmid 3 = q+1$.

\section{The conjecture}

Quoted verbatim from \texttt{conjectures/00000001006.md}:

\begin{quote}
\textbf{English.} Definition: An m-ovoid (meeting every generating line in
exactly m points). Conjecture: m-ovoids of Q(4,q) exist only when
m $\mid$ (q+1) (an existence criterion for m-ovoids).

\textbf{中文。} 定义：m-ovoid(每条生成线恰交 m 点)。猜想：Q(4, q) 的
m-ovoid 仅在 m $\mid$ (q+1) 存在(m-ovoid 存在判据)。
\end{quote}

The wording is identical in both languages, and the relation is written with the
divisibility symbol ``$\mid$''. We therefore refute the literal reading:

\begin{definition}[$m$-ovoid]\label{def:m-ovoid}
Let $q$ be a prime power and let $Q(4,q)$ be the parabolic quadric in
$\PG(4,q)$, regarded as a generalized quadrangle whose \emph{generating lines}
are its lines. An \emph{$m$-ovoid} is a set of points of $Q(4,q)$ that meets
every generating line in exactly $m$ points. A $1$-ovoid is simply called an
\emph{ovoid}.
\end{definition}

The claim to be refuted is thus the necessary condition
\[
\text{``an $m$-ovoid of $Q(4,q)$ exists''} \;\Longrightarrow\; m \mid (q+1).
\tag{$\ast$}
\]

\section{The quadric $Q(4,2)$}

For $q=2$ we use the standard vector model of $Q(4,2)$, which is the quadric in
$\PG(4,2)$ with quadratic form and polar form
\[
Q(x) \;=\; x_0 + x_1 x_2 + x_3 x_4, \qquad
B(a,b) \;=\; a_1 b_2 + a_2 b_1 + a_3 b_4 + a_4 b_3
\]
over $\F_2$ (here $x_i \in \F_2$ and $x_0^2 = x_0$). Thus
\[
Q(a+b) \;=\; Q(a) + Q(b) + B(a,b).
\]

\begin{definition}[points and lines of $Q(4,2)$]\label{def:q42}
The \emph{points} of $Q(4,2)$ are the nonzero vectors $x \in \F_2^5$ with
$Q(x)=0$. Its \emph{generating lines} are the triples
\[
\{a,\; b,\; a+b\}
\]
where $a,b$ are distinct points with $B(a,b)=0$ (equivalently, $a+b$ is again a
point).
\end{definition}

A direct enumeration gives the following well-known parameters.

\begin{proposition}\label{prop:params}
$Q(4,2)$ has exactly $15$ points and $15$ lines; every line contains exactly
$q+1 = 3$ points, and every point lies on exactly $q+1 = 3$ lines.
\end{proposition}

\begin{proof}
Enumeration over the $2^5 = 32$ vectors of $\F_2^5$. The number of nonzero
$Q$-isotropic vectors is $2^4 - 1 = 15$; each of the $15$ points is joined to
the other points $b$ with $B(a,b)=0$, giving $3$ lines per point and
$15 \cdot 3 / 3 = 15$ lines. See \texttt{reproduce.py} for the exact
computation and \texttt{lean4/Main.lean} for the formal verification.
\end{proof}

\section{An ovoid and its complement}

An ovoid of $Q(4,2)$ has $q^2+1 = 5$ points. The following is one.

\begin{theorem}\label{thm:main}
In the model of Definition~\ref{def:q42}, the set
\[
\ov \;=\; \bigl\{\, (0,1,0,0,0),\; (0,0,1,0,0),\; (1,1,1,1,0),\;
(1,1,1,0,1),\; (0,1,1,1,1) \,\bigr\}
\]
is an ovoid of $Q(4,2)$: it has $5$ points and meets each of the $15$ generating
lines in exactly $1$ point. Moreover:
\begin{enumerate}
\item every line of $Q(4,2)$ meets the complement $\ov^{c}$ in exactly
      $3 - 1 = 2$ points, so $\ov^{c}$ is a $2$-ovoid;
\item $|\ov^{c}| = 15 - 5 = 10 = q^2 + q$;
\item $m = 2$ does not divide $q+1 = 3$.
\end{enumerate}
Consequently the necessary condition $(\ast)$ is false, and with it conjecture
\texttt{00000001006}.
\end{theorem}

\begin{proof}
This is a finite verification. The five listed vectors are nonzero with
$Q=0$; enumerating the $15$ lines of Proposition~\ref{prop:params} gives exactly
one line through each pair of them (and no line containing two of them), so the
hit count is $1$ on every line. Since each line has exactly $3$ points and
$\ov$ occupies exactly one of them, the complement occupies exactly the
remaining $2$; the complement therefore consists of $15-5 = 10$ points and meets
every line in exactly $2$ points, i.e.\ it is a $2$-ovoid. Finally
$3 = 2 \cdot 1 + 1$, so $2 \nmid 3$. All counts are tabulated in
Section~\ref{sec:table} and checked exactly by \texttt{reproduce.py} and by
\texttt{lean4/Main.lean}.
\end{proof}

\begin{corollary}
Conjecture \texttt{00000001006} is false. The quadric $Q(4,2)$ possesses a
$2$-ovoid although $2 \nmid (2+1)$.
\end{corollary}

\begin{remark}[the general complement construction]
The argument is not special to $q=2$. If $Q(4,q)$ has an ovoid $\ov$, then
$|\ov| = q^2+1$ and the complement $\ov^{c}$ has $(q+1)(q^2+1) - (q^2+1) =
q(q^2+1)$ points; every line has $q+1$ points and meets $\ov$ once, hence meets
$\ov^{c}$ exactly $q$ times. So $\ov^{c}$ is a $q$-ovoid, while $q \nmid q+1$
for every $q \ge 2$. Ovoids of $Q(4,q)$ are known to exist for every prime
power $q$ (for $q=2$ the ovoid $\ov$ above is explicit), so the criterion
$(\ast)$ fails for \emph{every} $q \ge 2$ by taking $m = q$. This submission
formalises the smallest instance $q=2$; the general statement is recorded here
as context.
\end{remark}

\section{All $15$ lines and their hit counts}
\label{sec:table}

Each line of $Q(4,2)$ is listed as $\{a,b,a+b\}$ in coordinates
$(x_0,x_1,x_2,x_3,x_4)$; the last two columns give the number of points of the
line lying in $\ov$ and in $\ov^{c}$.

\begin{center}
\footnotesize
\begin{tabular}{llcc}
\toprule
\multicolumn{2}{c}{generating line} & $|\ell \cap \ov|$ & $|\ell \cap \ov^{c}|$ \\
\midrule
$\{(0,1,0,0,0),(0,0,0,1,0),(0,1,0,1,0)\}$ & & 1 & 2 \\
$\{(0,1,0,0,0),(0,0,0,0,1),(0,1,0,0,1)\}$ & & 1 & 2 \\
$\{(0,1,0,0,0),(1,0,0,1,1),(1,1,0,1,1)\}$ & & 1 & 2 \\
$\{(0,0,1,0,0),(0,0,0,1,0),(0,0,1,1,0)\}$ & & 1 & 2 \\
$\{(0,0,1,0,0),(0,0,0,0,1),(0,0,1,0,1)\}$ & & 1 & 2 \\
$\{(0,0,1,0,0),(1,0,0,1,1),(1,0,1,1,1)\}$ & & 1 & 2 \\
$\{(1,1,1,0,0),(0,0,0,1,0),(1,1,1,1,0)\}$ & & 1 & 2 \\
$\{(1,1,1,0,0),(0,0,0,0,1),(1,1,1,0,1)\}$ & & 1 & 2 \\
$\{(1,1,1,0,0),(1,0,0,1,1),(0,1,1,1,1)\}$ & & 1 & 2 \\
$\{(0,1,0,1,0),(0,0,1,0,1),(0,1,1,1,1)\}$ & & 1 & 2 \\
$\{(0,1,0,1,0),(1,1,1,0,1),(1,0,1,1,1)\}$ & & 1 & 2 \\
$\{(0,0,1,1,0),(0,1,0,0,1),(0,1,1,1,1)\}$ & & 1 & 2 \\
$\{(0,0,1,1,0),(1,1,1,0,1),(1,1,0,1,1)\}$ & & 1 & 2 \\
$\{(1,1,1,1,0),(0,1,0,0,1),(1,0,1,1,1)\}$ & & 1 & 2 \\
$\{(1,1,1,1,0),(0,0,1,0,1),(1,1,0,1,1)\}$ & & 1 & 2 \\
\midrule
\multicolumn{2}{r}{total over all lines} & $15$ & $30$ \\
\bottomrule
\end{tabular}
\end{center}

Every row has hit counts $1$ and $2$, so $\ov$ is an ovoid and $\ov^{c}$ is a
$2$-ovoid. (The $15$ lines above are exactly the whole line set; this
completeness is verified by enumeration in \texttt{reproduce.py} and by
\texttt{lines15\_is\_all\_lines} in \texttt{lean4/Main.lean}.)

\section{The divisibility failure}

The only arithmetic ingredient is
\[
q+1 = 3, \qquad m = 2, \qquad 3 = 2 \cdot 1 + 1 \;\Longrightarrow\; 2 \nmid 3 .
\]
So an $m$-ovoid exists for $m=2$, but $m$ fails to divide $q+1$. This
contradicts $(\ast)$ directly.

\section{Caveats and the exact reading}

\begin{itemize}
\item \textbf{The literal reading is the divisibility one.} Both the English and
      the Chinese statements write ``$m \mid (q+1)$'', i.e.\ ``$m$ divides
      $q+1$''. The counterexample above refutes precisely this reading.
\item \textbf{Weaker readings.} If a reader were to replace ``$m \mid (q+1)$''
      by ``$m \le q+1$'', then $m=2 \le 3$ would be compatible with the claim
      and the witness would not apply. But that is not what the filed statement
      says; the file contains no inequality and no other condition. We record
      this only to make the reading explicit.
\item \textbf{Only ``only when'' is used.} The statement is one-directional
      (a necessary condition), which is exactly what $(\ast)$ formalises; the
      counterexample violates that necessary condition, so it refutes the
      conjecture under either the one-directional or the biconditional reading.
\item \textbf{No algorithmic subtlety.} The ovoid $\ov$ and its complement are
      given by explicit vectors; every claim is a finite, exact computation
      with no search heuristics or probabilistic steps.
\item \textbf{Relation to the literature.} $Q(4,q)$ is a classical generalized
      quadrangle of order $(q,q)$; its ovoids (respectively, the $m$-ovoids it
      is known to admit) are a standard topic. Nothing in this note depends on
      any unpublished result: the $q=2$ instance is fully explicit.
\end{itemize}

\section{Formalisation and reproducibility}

The Lean~4 project under \texttt{lean4/} (core Lean only, \texttt{import Std},
no Mathlib, no \texttt{sorry}) models points as \texttt{Fin 5 → Bool}, defines
$Q$, $B$ and the line relation, hardcodes the ovoid $\ov$, and closes by
\texttt{decide}:
\begin{quote}\ttfamily
theorem O\_is\_ovoid : IsMOvoid 1 O = true\\
theorem Ocomp\_is\_2ovoid : IsMOvoid 2 Ocomp = true\\
theorem two\_not\_dvd\_three : \(\neg\) (2 \(\mid\) 3)\\
theorem two\_ovoid\_exists : \(\exists\) S, IsMOvoid 2 S = true\\
theorem lines15\_is\_all\_lines : allLinesE = lines15
\end{quote}
The axiom audit \texttt{lake env lean Check.lean} reports no \texttt{sorryAx}
(the finite facts depend at most on \texttt{propext}). The script
\texttt{reproduce.py} (Python~3 standard library only) rebuilds $Q(4,2)$,
enumerates all six ovoids, verifies the hit counts of $\ov$ and $\ov^{c}$, and
exits non-zero on any failure.

\begin{quote}\ttfamily
\$ python3 reproduce.py\\
... PASS: Q(4,2) has a 2-ovoid ... conjecture 00000001006 is FALSE.\\
\$ cd lean4 \&\& lake build \&\& lake env lean Check.lean\\
\$ tectonic main.tex --outdir build
\end{quote}

\end{document}
